UrbanFlow répond à six services (F1 à F6) par une API REST au-dessus de MongoDB.
Nous avons comparé trois déploiements sur les mêmes 20\,332\,093 lignes TLC de janvier
à juin 2024 : un serveur local, deux shards par zone de départ et deux shards par
période. Les 13 requêtes de la comparaison automatique donnent des réponses identiques
partout. Le placement par zone raccourcit tous les mélanges de E3 de 17 à 33\,\% par
rapport au local, et coûte 42\,\% de stockage en plus et, sur la charge large, 25\,\%
de CPU en plus. Le placement par période termine la charge large 7\,\% plus vite, mais
il est le plus lent sur les zones chargées et sur F5 à l'échelle de la ville (59 s contre
22 s). La perte d'un shard rend indisponibles les requêtes qui en ont besoin : l'API
renvoie alors une erreur 503 et ne présente jamais de résultat partiel.
